\subsection{Five-cohort bulk RNA-seq defines the established disease-state layer}\label{five-cohort-bulk-rna-seq-defines-the-established-disease-state-layer}

To characterize established disease tissue, we fit a pooled disease-versus-control model to fragment counts from five public liver RNA-seq cohorts (Suppli et al., 2019; Hoang et al., 2019; Govaere et al., 2020; Pantano et al., 2021; Chen et al., 2023). The unit of inference is the genetic-identity cluster, with one library per cluster. These clusters are not verified independent donors, but no two selected libraries, within or across cohorts, had genotypes consistent with one person (Methods). The model includes 774 clusters (138 control and 636 disease; 326, 207, 140, 71 and 30 in GSE213621, GSE135251, GSE162694, GSE130970 and GSE126848), tests 24,115 Ensembl gene rows and has fixed cohort and inferred-sex terms (Figures 3A and 3B). It gave 1,854 canonical calls at Benjamini-Hochberg (BH) padj \textless{} 0.05 across all 24,115 genes and \textbar log2FC\textbar{} \textgreater{} 0.5. These are tissue-level differences, so a call can reflect a change in cell-type proportions, a change in expression within cells, or both. We keep colocalization and RNA measurements as separate evidence throughout.

To define what a missing call means, we permuted disease labels within cohort and then added known log2 fold changes to the 774-unit counts (50 replicates). The pooled fit called a uniform 1.0 log2FC change in 50.7\% to 100.0\% of genes per expression decile and called 0.0\% of unchanged genes. We call the smallest change detected with 80\% recall the detection limit. It was 0.649 log2FC in decile 8, the lowest of any decile, 0.653 in the most highly expressed decile and above 1.5 in the least expressed decile. In the two lowest expression deciles, a missing call therefore does not exclude a 2-fold change. Those genes are below the detection limit, not unchanged.

We next fit each cohort separately with the canonical threshold (Figure 3C). Four cohorts use the 24,115-gene family of the pooled fit. Because GSE126848 calls track sequencing depth (see below), we refit that cohort on the 15,592 genes that pass its own expression filter (filterByExpr), with BH correction over those 15,592 genes. The refit gives 3,440 calls. The union of per-cohort calls contains 12,180 genes: 8,341 called in one cohort and 3,839 in two or more. Only 2,628 of the 3,839 have the same sign in every cohort that calls them (1,820 in two cohorts, 582 in three, 186 in four and 40 in all five). The other 1,211 (31.5\%) are called up in one cohort and down in another, and 946 of these include a GSE126848 call. Without GSE126848, 346 of 2,565 multi-cohort genes (13.5\%) have mixed signs. GSE135251 is one of the cohorts for 325 of these 346 genes, and without it as well the share is 2.2\%. With the 6,147 GSE126848 calls from the 24,115-gene family in place of the refit, 1,865 of 4,696 multi-cohort genes (39.7\%) have mixed signs. These 17,347 cohort memberships are a different set from the 1,854 pooled calls, and shared membership is not replication. In 50 random splits of the 493 SteatoSITE donors into five groups of about 99, 26.8\% of called genes (2.5-97.5\% range 24.8 to 29.9\%) were called in only one group, and 0.17\% of genes called in two or more groups had mixed signs (F3--F4 versus F0--F1, BH q \textless{} 0.05 within each group). This within-cohort stage comparison produces less disagreement than the five disease-versus-control cohorts. Because the contrasts, group sizes, protocol and case mix differ, it does not partition the causes of between-study disagreement. GSE135251 membership cannot be matched for genetic ancestry, because only 2 of its 10 controls are of European genetic similarity and 6 are of African similarity.

\begin{mainfigureplate}{3}{654.48bp}
\noindent\makebox[\linewidth][c]{\paperpanel{A}{216bp}{figures/main/figure3/fig3a_cohort_metadata_matrix.pdf}\hspace{8bp}\paperpanel{B}{130bp}{figures/main/figure3/fig3b_nas_fib_grid.pdf}\hspace{8bp}\paperpanel{C}{230bp}{figures/main/figure3/fig3c_cohort_alluvial.pdf}}
\par\vspace{8pt}
\noindent\makebox[\linewidth][c]{\paperpanel{D}{180bp}{figures/main/figure3/fig3d_pca_fibrosis_gradient.pdf}\hspace{8bp}\paperpanel{E}{200bp}{figures/main/figure3/fig3e_stage_remodeling.pdf}\hspace{8bp}\paperpanel{F}{170bp}{figures/main/figure3/fig3f_steatosite_f3f4_scatter.pdf}}
\par\vspace{8pt}
\noindent\makebox[\linewidth][c]{\begin{minipage}[t]{240bp}\vspace{0pt}\paperpanel{G}{240bp}{figures/main/figure3/fig3g_steatosite_composition.pdf}\par\vspace{8pt}
\paperpanel{H}{240bp}{figures/main/figure3/fig3h_within_hepatocyte_partition.pdf}\end{minipage}\hspace{8bp}\paperpanel{I}{258.48bp}{figures/main/figure3/fig3f_stage_deg_genetics_matrix.pdf}}

\end{mainfigureplate}
\papercaption{3}


GSE126848 calls track sequencing depth as well as disease. Its model on the fixed 24,115-gene family contains 23 control and 7 disease clusters, all male, and gives 6,147 calls. In that historical fixed-family fit, direction agrees with the pooled fit for 84.36\% of the 1,854 pooled calls. Under the current native filter, 937 of 1,056 testable pooled calls agree (88.7\%; Figure S3AD); the other 798 genes are untestable in this cohort. The source identity screen removed 24 of 31 disease records, including all ten women, and 2 of 25 controls. Restoring those ten female disease libraries in a sex-adjusted sensitivity model ($n = 40$; Figure S3AD) yielded 1,803 testable pooled calls (84.9\% same sign, 1,531 of 1,803; 38.7\% called same sign; all-gene Spearman rho = 0.9998 against the male-only fit). This unequal eligibility may limit how representative the cohort is. The retained disease libraries also have about half the fragment count of the controls on the tested genes (median 19.5 versus 35.3 million). Because voom adds 0.5 to every count before computing log-CPM (Law et al., 2014), a gene with no reads has a higher log-CPM in shallower libraries. Of the 6,147 calls, 437 are such genes, with zero reads in all 30 libraries. More broadly, 2,730 of the 6,147 calls are on genes that the cohort's own expression filter removes, and 2,639 of these are higher in disease. Adding log2 library size as a covariate keeps 2,025 to 3,056 of the 6,147 calls, depending on whether library size is counted before or after normalization, and keeps none of the 437 (Supplementary Note 7). We therefore do not read GSE126848 per-cohort calls as disease replication.

Historical direction agreement among the 1,854 pooled calls ranged from 78.3\% to 100\% across the five cohorts, including the unfiltered GSE126848 comparison (Figure S3C). When we removed each cohort and refit the full model, 66.7\% to 96.8\% of the 1,854 calls were still called with the same sign, and among calls still tested, 99.7\% to 100\% kept their sign. Removing GSE213621 lost the most calls (618, of which 108 were no longer tested). Fitted alone, GSE213621 calls 1,654 of the 1,854 pooled calls (89\%) with the same sign, and the next highest, GSE162694, calls 1,014 with the same sign. Leave-out fits on a fixed gene set gave all-gene Spearman correlations of 0.870 to 0.995 and threshold retention of 70.7\% to 96.6\% (Figure S3C). Neither procedure is independent replication of a pooled result. In a spike-in calibration on these counts, an effect placed only in GSE213621, which carries 45.1\% of the summed quality weight, was called in the pooled fit in 87.45\% of replicates. An effect placed in any other single cohort was called in at most 0.25\% of replicates (1.0 log2FC in both cases). Keeping a call in all five leave-one-cohort-out refits therefore does not show that a real gene's call is independent of any single cohort. Adjusting for log library size, top-100-gene share and detected genes together kept 1,563 (84.3\%) or 1,569 (84.6\%) of the 1,854 calls under two library-size definitions. Across these two fits and four fits that add one library property each, all-gene log2FC Spearman correlations with the primary fit were 0.979 to 0.992 (Supplementary Note 7). Whether these library properties are technical or reflect disease biology is not resolved. GSE213621 carries a median 54\% of the per-gene information in the pooled fit, and 409 of the 1,854 calls are called in GSE213621 alone among the four cohorts other than GSE126848. In GSE213621, type 2 diabetes is present in 47.8\% of NAFLD participants and 17.4\% of controls (Chen et al.~2023, PMID 36535507), and diabetes is not in the model. The pooled coefficient is an information-weighted, cohort-adjusted tissue association; it does not establish that the effect transfers to each cohort. We therefore evaluated the complete 24,115-gene family with restricted-maximum-likelihood random effects and modified Hartung--Knapp uncertainty, using each cohort's current expression eligibility. Of the 1,854 fixed pooled calls, 1,850 random-effects means retained the pooled sign and 1,208 had the same sign in every available cohort. The smallest full-family BH q was 0.1714, so none passed this sensitivity. Median I-squared among the fixed calls was 60.35\%. GSE126848 contributed estimates for 1,056 of the 1,854 genes; its other 798 entries were untestable, not zero effects. This analysis has only four or five cohorts per gene and uses moderated input standard errors. It therefore establishes uncertainty about between-cohort mean effects, not their absence. Earlier DerSimonian--Laird classifications used different cohort inputs and uncertainty rules and remain historical descriptions (Supplementary Note 19). Neither those classes nor agreement within discovery cohorts establishes independent disease-versus-control replication.

Within the same cohorts, controls include more units of African genetic similarity than disease biopsies do (31 of 138 against 37 of 636). We added four genetic ancestry principal components (PCs), computed from RNA-called genotypes, to the disease-versus-control model. This changed 233 of the 1,854 calls (140 lost, 93 gained), more than within-cohort-and-group permutations of the same PCs (Holm-adjusted \emph{p} = 0.0199 over two tests). Restricting the fit to units called as European genetic similarity (EUR) changed 570 calls (Holm-adjusted \emph{p} = 0.0199) (Supplementary Note 7). In each test, one reference fit changed at least as many calls as observed (one of 200 permuted-PC fits and one of 100 random subsets), so the precision of both p values is limited by the number of reference fits. The changes are threshold crossings. No lost call changed sign, and 127 of the 140 lost calls fell to \textbar log2FC\textbar{} at or below 0.5 with padj still below 0.05. Over the 1,854 calls, log2FC with the PCs correlates with the primary values at Pearson 0.9972 (descriptive). The call count is therefore sensitive to ancestry near the 0.5 threshold. In total, 64 genes, GSTM1 among them, met the ancestry-sensitive rule. This label is descriptive, and the primary fit and its 1,854 calls are unchanged.

To test whether disease-state genes also vary with liver fat outside biopsy cohorts, we used GTEx liver (GTEx Consortium, 2020). In GTEx, fat was measured on each donor's liver H\&E slide as the clear-vacuole area fraction (Supplementary Note 19). We tested 1,002 disease-state calls measurable in GTEx liver, in 222 postmortem donors not selected for liver disease. Across these genes, changes with image-measured histologic fat went in the same direction as the disease-versus-control changes in biopsies (rho 0.638). Of the 326 genes with fat q \textless{} 0.05, 323 have the same sign as in the pooled biopsy fit. This held in the 138 donors who did not die on a ventilator (rho 0.632). The genes also followed fat more than ischemic time (rho 0.636 for fat and -0.350 for ischemic time in a joint model). Two negative controls did not meet their predictions. Slide tissue area gave rho 0.354 (\emph{p} = 0.056; predicted \textbar rho\textbar{} \textless{} 0.2), and ischemic time gave rho -0.350 (predicted -0.2 to 0.3). The prespecified stop rule for tissue area required both \emph{p} \textless{} 0.05 and \textbar rho\textbar{} of at least half the fat rho (0.319). Only the second condition held, so the rule did not stop the analysis. Both positive controls met their predictions. PLIN2, which is not a disease-state call, rose by 1.07 log2 CPM from the lowest- to the highest-fat donor (moderated \emph{p} = 7.5 × 10⁻⁷). Across 197 steatosis-associated genes that are not disease-state calls, the GTEx fat slope correlated with the biopsy steatosis effect (rho 0.717, Freedman-Lane \emph{p} = 0.0002, the smallest two-sided value 10,000 permutations allow). In a post hoc sensitivity, adding slide tissue area to the model left the fat correlation almost unchanged (rho 0.631 against 0.638; difference -0.0065, paired donor-bootstrap 95\% interval -0.033 to 0.004). The fat measured here is histologic steatosis of unrecorded cause in postmortem liver. This result does not assign any biopsy-cohort call to steatosis.

We then scored the adult signature in two pediatric cohorts (Supplementary Note 19). The adult signature score is built from the 1,854 adult disease-versus-control calls, and 1,064 of their gene symbols map to the GSE66676 array. GSE66676 profiled liver from severely obese adolescents at bariatric surgery, so its controls are also severely obese (Xanthakos et al., 2015). In GSE66676, the score did not differ detectably between 32 adolescents with NAFLD and 33 without (standardized difference d 0.44, 95\% CI -0.07 to 0.95). The cohort had 80\% power for d of 0.71, and the median adult NAFL-versus-control shift (d 0.97) lies outside this interval. In 52 children in GSE185051 (Yao et al., 2022), the score correlated with NAS at rho 0.12 (nominal \emph{p} = 0.41).

We then tested the adult calls in two adult cohorts outside the pooled fit (Supplementary Note 19). GSE48452 (Ahrens et al., 2013) profiled liver on microarrays from 12 lean participants, 16 obese participants with normal histology, 9 with steatosis and 17 with NASH. Over 1,036 calls matched on the array, effects against lean controls scaled with the discovery effects for NASH (slope 0.336; label-permutation Holm \emph{p} = 0.003 over three contrasts; sign agreement 0.84) but not for obese normal liver (slope −0.006, Holm \emph{p} = 0.56). The obese slope was −0.02 of the NASH slope (participant bootstrap 95\% interval −0.61 to 0.31), so the calls are not explained by obesity alone in this cohort, although up to about a third of the effect could be shared with it. In GSE281797 (Zhao et al., 2026), 62 participants with MASL or MASH and 29 controls, 1,063 calls were testable. Sign agreement was 0.653, which did not exceed the label-permutation null (95th percentile 0.70; Holm \emph{p} = 0.13 over two statistics). Effects scaled at about one fifth of the discovery size (slope 0.209, Holm \emph{p} = 0.016), and no single call reached BH q \textless{} 0.05 in the replication direction. With 29 controls, the cohort has little power per gene. In a post hoc breakdown, calls in the consistent class agreed in sign more often (0.72; 490 calls) than cohort-discordant calls (0.50; 48), with no test.

We next asked whether a prespecified phenotypic endpoint, rather than groups chosen to improve unsupervised separation, reproduces the disease-state comparison (Figure 3D). The endpoint compares 113 strict controls with 279 advanced-disease clusters across four cohorts. Advanced disease means NAS of at least 5 or fibrosis of at least F3. Strict controls have no observed NAS above 1 and no observed fibrosis above F0. Missing histology is allowed only when no observation contradicts source control status. Each cohort must provide at least five units in each arm. Expression is adjusted for cohort and inferred sex without protecting endpoint labels, and the density and scatter plots show the same 392 clusters. Across the complete 24,115-gene family, endpoint effects correlate with pooled effects at rho = 0.953. Under the same threshold, applied independently, the endpoint recovers 1,755 of 1,854 canonical calls and makes 1,195 endpoint-only calls. This is a related sensitivity analysis in overlapping units. It does not replace the five-cohort disease-versus-control estimand.

Multicohort differential expression and stage-ordered transcriptomic summaries in MASLD are established analyses, not contributions of this paper. A nine-cohort public bulk meta-analysis of 897 NAS-scored and 856 fibrosis-scored biopsies reports stage-resolved differential expression and pathway change on overlapping data, including three of the five cohorts used here (Kumak et al., 2026). It follows earlier harmonized liver resources (Piras and DiStefano, 2024; Alvarez-Sola et al., 2026). We therefore claim no novelty for recovering a multicohort disease-state gene set or for ordering expression change by histologic stage. What Figure 3 adds is the estimand and the unit: one fragment-count dataset, genetic-identity clusters as the inferential unit in every bulk fit, a complete 24,115-gene family carried through each contrast, and the connection of these changes to the inherited map (Figures 3F and 3G). The five cohorts include no repeat-biopsy cohort. Where the Resource uses repeat biopsies, they enter only within-person paired models and never a cross-sectional contrast.

\subsection{Expression differs from F0 at F2 to F4, and the F1 difference depends mainly on one cohort}\label{expression-differs-from-f0-at-f2-to-f4-and-the-f1-difference-depends-mainly-on-one-cohort}

To describe cross-sectional stage-associated remodeling, we compared each fibrosis stage from F1 to F4 with a common F0 reference (Figure 3E). We used the three cohorts that record an exact F0--F4 stage (GSE130970, GSE135251 and GSE162694; 387 staged units, F0/F1/F2/F3/F4 = 103/102/88/69/25). GSE213621 records only F0--F1, F2 and F3--F4 bins (with source Controls set to unstated fibrosis stage, leaving 258 staged disease units in GSE213621 and 367 staged evaluation units across GSE162694 and GSE213621) and is not used in these exact-stage contrasts. Each model has fixed cohort and inferred-sex terms when estimable, with a pre-voom per-cohort filter requiring at least three positive-count samples (24,113, 24,112, 24,111 and 24,109 genes tested across F1--F4) to prevent zero-read depth artifacts. BH correction covers each full tested family, and there is no stage-specific fold-change filter, so these counts are distinct from the canonical disease-versus-control calls. The contrasts gave 1,338, 4,269, 6,190 and 7,042 genes at BH q \textless{} 0.05 for F1, F2, F3 and F4 versus F0 (182, 798, 1,756 and 4,018 also with \textbar log2FC\textbar{} \textgreater{} 0.5). These counts depend on sample size, effect and uncertainty, and they cannot locate where remodeling begins. Figure 3E also shows continuous F0-centered activity for the five available axes of a six-factor non-negative matrix factorization (NMF) of bulk expression (k6 NMF). These axes describe continuous variation, not patient classes. The sixth axis, k6-P6, is not shown. With its fixed weights restricted to the 4,517 retained genes, its scores recomputed from the original counts correlated with the released scores at Spearman 0.9874, below the prespecified minimum of 0.99. On the fragment counts used here, the five shown axes correlate with the released scores at Spearman 0.928 (k6-P5) to 0.998 (k6-P1) over 617 shared libraries. Full fibrosis, NAS and composition views are in Figure S3D--L.

The F1 count depends mainly on GSE130970. Without it, 14 of the 1,338 genes remain (155 units). Fitted alone, the cohorts give 252 (GSE130970, 50 units), 3 (GSE135251, 90 units) and 0 (GSE162694, 65 units) genes (Figure S3F). In a cohort-by-stage model, 95.8\% of the 1,338 effects have the pooled sign in all three cohorts. The cohorts therefore mostly agree on the direction of the F0-to-F1 differences, but only GSE130970 gives many calls on its own. With F0 restricted to its 89 disease biopsies (191 units), F1 versus F0 gives 783 genes (Figure S3E). The F2, F3 and F4 versus F0 counts stay between 1,934 and 9,546 when any one cohort is omitted. The later adjacent steps do depend on single cohorts. F1 to F2 gives 0 genes without GSE162694 (133 units), and F2 to F3 gives 2 without GSE135251 (54 units). Effect sizes also differ between cohorts at F4, where 4,047 of the 7,043 F4-versus-F0 genes have a cohort-by-stage interaction at BH q \textless{} 0.05. We prespecified a test of whether within-cohort library size, library Gini coefficient and normalization factor account for the fibrosis- and NAS-stage calls. We could not interpret that test, because its positive control, a planted covariate, failed. The Figure 3E counts are unchanged. As a post hoc description, Supplementary Note 7 gives the calls in each stage contrast that stay significant with these properties added.

We also asked whether grouped fibrosis-stage effects from these three cohorts appear in GSE213621 (Supplementary Note 7). This check is not blind. GSE213621 stage labels already enter the continuum analysis. We had also read a fit of the same F3--F4 versus F0--F1 contrast in GSE213621, and its sign agreement with the discovery cohorts, before the analysis plan was written. The result therefore describes the size of the agreement and its interval. It is not a prediction confirmed on unseen data. We refit the three cohorts on the same bins in disease units only (371 units), with controls excluded on both sides. Grouped effects point the same way in the 258 GSE213621 disease units. The slope was b = 0.708 (95\% bootstrap interval 0.550 to 0.867) over the 8,449 genes that have BH q \textless{} 0.05 for F3--F4 versus F0--F1 in the three-cohort fit and pass the GSE213621 expression filter. This is larger than the slope that submission order and sex structure produce (permutation null mean 0.006; Holm-adjusted \emph{p} = 0.002). Whether this agreement is as large as the agreement among the discovery cohorts is not determined (ratio 0.793, 95\% interval 0.488 to 1.098). The F2 versus F0--F1 contrast was untestable, so we claim nothing about F2 (Supplementary Note 7). A negative control failed. The GSE213621 male-versus-female coefficient is also associated with the discovery stage effects (slope -0.155, 95\% interval -0.279 to -0.031, which excludes zero). Either sex and fibrosis bin are not separated in GSE213621, or sex and fibrosis share real biology, for example in sex-biased hepatocyte metabolic genes. Figures 3E and 3F remain three-cohort results.

We then tested the same stage effects in a cohort that no part of this Resource had used. SteatoSITE (Kendall et al., 2023) deposited total RNA from archival FFPE liver with NASH-CRN fibrosis stage (ENA PRJEB58625). We wrote the analysis plan, with its families, controls and pass rules, before reading any SteatoSITE expression. After blinded quality control and a genotype identity screen, 493 staged donors remained (source stage code 1, read as F0--F1, 255; F2 74; F3 75; F4 89). The stage-expression and NMF models include sequencing batch and inferred sex; the exon/intron-ratio model is specified separately in Methods. In F3--F4 versus F0--F1, 8,460 of the 8,560 discovery genes were measurable. 93.5\% had the discovery sign, against 52.6\% under a stage permutation within batch and sex. The slope was b = 1.462 (95\% CI 1.323 to 1.601; Holm-adjusted p = 0.002 over the two grouped contrasts; Figure 3F; per-contrast summary in Figure S3Y). The prespecified positive controls passed, and unlike in GSE213621 the sex negative control passed (slope -0.052, 95\% CI -0.159 to 0.054). A larger sample improves precision but does not itself explain a slope above 1; stage mixture, tissue composition and protocol differ between the sources. The F2 versus F0--F1 contrast and the three adjacent steps of Figure 3E agreed in sign for 94\% to 96\% of genes, but each was untestable under the prespecified power rule (minimum detectable slope 0.36 to 0.59), so we read them as descriptive. GNMT, MAT1A and CYP2C19 fall with stage in SteatoSITE (Holm over three, all p \textless{} 10\^{}-35), and 28 of the 29 later-step Figure 3I marks replicate (BH within 29). Stage-associated expression in SteatoSITE, including these genes, was reported by Kendall et al.; what we add is the test of our own fixed effects. Adding seven isometric log-ratio cell-composition coordinates from a 22-autosome leave-one-chromosome-out (LOCO) deconvolution that separates hepatocytes from cholangiocytes reduced the F3--F4 slope to 0.224 (60.1\% same sign; 1,132 genes at BH FDR \textless{} 0.05, compared with 1,174 under non-LOCO 8-part deconvolution, Spearman rho = 0.9973; Figure 3G), so most of the bulk stage signal cannot be separated from tissue composition in bulk data. This does not mean that most of it is a change in cell mix. Even with LOCO chromosome exclusion preventing single-gene mathematical circularity, the deconvolved fractions are estimated from the same libraries and move with stage. Of the 7,214 replicated stage genes, 1,276 changed in the same direction inside hepatocytes in the uncorrected single-nucleus partition (C3; Figure 3H); after empirical ambient-RNA subtraction (C3b, 0 of 9 non-hepatocyte negative controls called, donor-label permutation \emph{p} = 0.000999 with ambient correction refit inside every permutation), 188 genes remained supported within hepatocytes (31 higher and 157 lower with stage; Figure 3H). Because these libraries keep intronic reads, we could also compare intronic and exonic stage effects using joint voom inverse-variance weights. The change in intronic reads, which mark pre-mRNA, followed the change in exonic reads across 12,010 genes (Spearman 0.949; 2,727 exon-minus-intron ratio calls at q \textless{} 0.05, 96 with \textbar $\Delta$E $-$ $\Delta$I\textbar{} $\ge$ 0.5, batch-by-sex permutation \emph{p} = 0.005; Figure S3AB). For GNMT, MAT1A, COL1A1, IGFBP7 and TNFSF10 the intronic change was 90\% to 113\% of the exonic change, consistent with a substantial transcription-associated tissue component. This comparison cannot separate changes in cell mix, transcription, RNA processing or stability. Four of the Figure 3E NMF axes (P1 to P4) also showed the discovery stage pattern in SteatoSITE (one-sided permutation tests, Holm over five axes, each p = 0.005; Figure S3AC). These are cross-sectional differences, and FFPE total RNA differs from the discovery protocols, so the comparison uses signs and slopes, not expression levels (Supplementary Note 17).

We kept histology as the clinical severity anchor and asked a separate question: whether an externally anchored transcriptomic continuum could resolve molecular variation that fibrosis categories compress (Figures S3I--S3U). All five cohorts were scored and enter descriptive continuum-window displays, and the four cohorts with recorded fibrosis enter stage displays. We evaluated two prespecified co-primary scores inferentially in GSE162694 and GSE213621, with fibrosis stage and inferred sex in every genetic-cluster-level model. The continuum recovered the broad stage-associated direction of remodeling in transcripts, pathways, NMF axes and Hotspot programs (117 prespecified single-cell gene programs, each a weighted set of co-expressed genes; Figure 4) and kept signal after adjustment for stage (Figures S3M--S3U and S4H--S4N). For the Hotspot programs, that signal depends on cell-type composition adjustment. In the 367 disease units, without the 68 GSE213621 controls, 94 of the 116 estimable programs (of 117) are continuum-associated (BH q \textless{} 0.05) without adjustment and 53 to 76 after adjustment, depending on the adjustment arm and the deconvolution input. These comparisons use the same expression data, so they are alignment analyses, not independent validation. Sliding windows are fixed descriptive summaries only. This extension does not change Figure 3, the definition of the canonical disease-versus-control calls or the interpretation of NMF axes.

\subsection{GNMT, MAT1A and CYP2C19 have lower RNA and protein abundance in disease liver}\label{gnmt-mat1a-and-cyp2c19-have-lower-rna-and-protein-abundance-in-disease-liver}

GNMT, MAT1A and CYP2C19 are three hepatocyte metabolic genes that illustrate the disease-state result. All three are lower in the pooled disease-state comparison and in adjusted liver protein (Figure S3X). Lower GNMT and MAT1A in MASLD liver have been reported before (Fisher et al., 2009; Murphy et al., 2013; Zubiete-Franco et al., 2016; Fernández-Tussy et al., 2019; Männistö et al., 2019), so these results reproduce reported decreases. The F0 group contains 14 source controls and 89 disease biopsies, so we also refit the stage contrasts with a disease-only F0 and with an added control-status term. All three genes have lower abundance at F2, F3 and F4 than at F0, with BH q \textless{} 0.001 under every F0 reference and without GSE130970. Omitting GSE135251 or GSE162694 leaves every effect negative. Every q is ≤ 0.034 except MAT1A at F4 without GSE135251 (q = 0.059; 11 F4 units). The differences at earlier steps depend on single cohorts (Supplementary Note 7), and we do not use them to place where the lower abundance begins. MAT1A falls below the canonical pooled 0.50 effect floor. It passes the interval-null sensitivity, a test that the absolute log2FC exceeds 0.25 (BH q \textless{} 0.05). We do not relabel it as a canonical call. Lower bulk abundance of these genes can reflect a smaller hepatocyte share of the biopsy as well as lower per-cell abundance. We compared donor hepatocyte pseudobulk from healthy donors with that from steatohepatitis or cirrhosis in four single-nucleus cohorts. GNMT and MAT1A were lower in GSE298719 and GSE256398 (GNMT -2.10 and -1.85 log2 CPM), and CYP2C19 was lower in GSE256398 (the GSE298719 probe panel lacks it). In the two cohorts not used to find this (GSE202379 and GSE244832; 8 healthy and 44 disease deduplicated donors), the prespecified HC3-studentized combined test was uninformative: GNMT +0.24 (permutation \emph{p} = 0.750, Holm \emph{p} = 0.750), MAT1A -0.37 (\emph{p} = 0.313, Holm \emph{p} = 0.691) and CYP2C19 -1.06 (\emph{p} = 0.230, Holm \emph{p} = 0.691). With 3 and 5 healthy donors, these two cohorts could not detect GNMT or MAT1A changes of the size seen in the first two, and the GNMT effect differed between cohorts (Supplementary Note 19). We then applied the same HC3-studentized Freedman--Lane model to GSE289173 (Figure 4E), a cohort not used in any earlier analysis of these genes, with 23 control and 17 steatohepatitis donors (excluding one unconfirmed donor, GSM8785704). All three genes were lower per hepatocyte: CYP2C19 -1.69 log2 CPM (95\% CI -2.56 to -0.81; Holm \emph{p} = 0.0034), GNMT -1.17 (-1.92 to -0.41; Holm \emph{p} = 0.0045) and MAT1A -0.92 (-1.44 to -0.41; Holm \emph{p} = 0.0034). Each decrease was larger than 80\% of decreases among genes with the same control expression, so it is not explained by general dilution of hepatocyte RNA. In this cohort control and disease livers were collected as separate series, so tissue source cannot be separated from disease. Lower per-hepatocyte abundance is therefore seen in four of five cohorts for GNMT and in all four cohorts that measured CYP2C19. These data support a recurring metabolic abundance pattern. They do not show a decline over time within a person, a shared upstream regulator or a patient subtype.

Within 46 MASLD participants with liver proteomics (Boel et al., 2025), the three proteins also correlate positively with each other after adjustment for acquisition batch, age, BMI and sex (rank-residual partial Spearman rho = 0.406--0.472 over the three pairs; three-pair BH q = 0.00243--0.00772, accounting for nuisance-model degrees of freedom; Figure S3X). All adjustment comparisons use the same 46 participants. Adding NAS reduces the correlations to 0.223--0.277. Paired participant-bootstrap intervals for these reductions exclude zero, but neither the reduction nor a failed significance threshold establishes mediation (Supplementary Note 7). This example is selection-conditioned, because the protein cohort contributed to its selection. It supports a recurring metabolic abundance pattern, and independent liver-protein confirmation is unresolved. The covariation is also not specific to these three proteins. Random triplets of proteins that are lower in disease reach the trio's mean rho in 5.1\% of draws.

\subsection{No association was detected between colocalization support and disease-state RNA change across 3,587 genes}\label{no-association-was-detected-between-colocalization-support-and-disease-state-rna-change-across-3589-genes}

Figure 3I places stage-associated expression beside regulatory genetics and keeps the two measurements separate. Fifty-four genes belong to the 430-gene primary SuSiE set after the PSD eQTL refit and pass BH FDR \textless{} 0.05 in at least one adjacent fibrosis contrast. The adjacent contrasts (F0 to F1, F1 to F2, F2 to F3 and F3 to F4) contain 103/102, 102/88, 88/69 and 69/25 units (205, 190, 157 and 94 in total), and no identity cluster enters both sides of a contrast. Each contrast is its own BH family, so the union carries no joint 5\% FDR. When the 5\% budget is split across the four contrasts (BH q \textless{} 0.0125 in each), 15 genes pass (16 filled dots). This bounds the FDR over all discoveries in the four tested families at 5\%, but not specifically within the SuSiE subset. Across the four contrasts, the 54 genes give 216 displayed gene-contrast effects, of which 59 are significant in their complete BH family (25 selected at F1--F2 or later and 29 selected only by F0--F1). Twenty-nine of the 54 genes are selected only through an F0-to-F1 mark, and those marks share the cohort dependence of the F1 count. None of the F0-to-F1 marks keeps BH q \textless{} 0.05 without GSE130970, and 18 keep it with disease-only F0. Later marks also depend on single cohorts. None of the 13 F1-to-F2 marks keeps q \textless{} 0.05 without GSE162694, and 1 of the 59 marks (OMG, F2 to F3) keeps it after each single-cohort omission. Each omission also removes 12\% to 69\% of a contrast's units, which lowers power. For each gene, the genetics column shows the PP.H4 of one exact gene-trait pair: the gene with its strongest eligible driving GWAS trait. Fine-mapping PIPs are not shown. Gene selection is conditioned on stage significance, so the panel establishes neither enrichment nor causality. Each gene keeps its evidence state, linkage alternative and next discriminating experiment.

Colocalization does not place GNMT, MAT1A or CYP2C19 in an inherited regulatory program. None of the three has a primary SuSiE PP.H4 above 0.5 (GNMT 0.128, from the UK Biobank ALT GWAS). As described for Figure 2, GNMT's only colocalization above 0.5 is a single-signal ABF PP.H4 of 0.706 in a European NAFLD GWAS (Ghodsian et al., 2021). Multi-signal SuSiE colocalization returned no result for GNMT in that study, as for 18,939 of the 18,961 gene rows in that study. There is therefore no multi-signal estimate to confirm or contradict the single-signal result. GNMT is a provisional link between the genetic and disease-state maps. The RNA, stage and protein observations stand independently of that result.

We plotted colocalization support against disease-state log2FC in the 3,587 genes evaluated under the adopted priority eligibility rule and the PSD eQTL refit (Figure S3Z). Because PP.H4 is unsigned, the test uses the magnitude of RNA change. Spearman rho between PP.H4 and \textbar log2FC\textbar{} was 0.00094 (nominal \emph{p} = 0.955). The 420 supported genes split 216 up and 204 down. The displayed groups Neither, Transcriptomics only, Genetics only and Both contain 2,984, 183, 394 and 26 genes. Without expression stratification, 26 shared genes were observed against 24.5 expected (Fisher OR 1.076, nominal \emph{p} = 0.739). The expression-decile test gives 20.9 expected (95\% null range 13 to 29; Mantel-Haenszel OR 1.321, 95\% CI 0.851 to 2.051; exact \emph{p} = 0.222; Figure S1A). These values follow the shared-posterior check of carried negatives and the PSD eQTL refit; the earlier 3,769-gene panel is historical. Under the alternative any-eligible-study rule, 3,734 genes are evaluated (420 genetic, 217 RNA, 26 shared against 20.6 expected; MH OR 1.340, 95\% CI 0.864 to 2.080; exact \emph{p} = 0.220); this is a sensitivity rather than the adopted rule. The native RNA interface and the annotation-adjudicated Catalog have different denominators, and one does not replace the other.

We separately compared 426 one-to-one RNA-abundance-matched pairs, each a genetic gene and a control gene, using the same interval-null sensitivity (lfc = 0.25, BH q \textless{} 0.05) (Figure S3G). Calls occur only in the genetic member in 37 pairs, only in its control in 28, in neither in 360 and in both in one (paired OR = 1.321; nominal exact McNemar \emph{p} = 0.321). Genes are not biological replicates. Because PP.H4 is unsigned, neither comparison tests risk-allele direction or establishes biological opposition or complementarity.

The inset of Figure S3Z compares lncRNA with protein-coding genes in the 774-unit RNA fit. Over 21,501 genes, the odds of a call for lncRNA relative to protein-coding genes were 1.12 without adjustment (95\% CI 1.00--1.25) and 0.43 after adjustment for expression and detection (95\% CI 0.38--0.49). Calls also differed in direction: 83.0\% of 1,157 protein-coding calls and 47.9\% of 549 lncRNA calls were up. We matched 520 one-to-one lncRNA/protein-coding pairs on covariates at an outcome-blind 0.15-SD caliper. In these pairs, the conditional odds ratio for upregulation, lncRNA versus protein-coding, was 0.132 (30 discordant pairs with the lncRNA member up against 227 with the protein-coding member up; Wald 95\% CI 0.090--0.193). Gene-level matching is descriptive, not biological replication.

Having described the bulk disease-state and stage layers, we next examined prespecified cell-type-specific gene programs and regulatory activities at the level of biological donors (Figure 4).
